\section{Attribute Synthesis: Algorithm and Guarantees}
\label{sec:attribute-synthesis-algorithm-and-guarantees}
We now synthesize RHS attributes from the LHS axis representation in three
passes: determine surviving axes, check ordering feasibility, and construct
attributes. We call this algorithm the \emph{complete constructor}: within
Section~\ref{sec:formal-scope-and-guarantees}'s domain, it finds valid RHS axis attributes
whenever they exist.

\subsection{Canonical Intermediate Axes}
\label{sec:canonical-intermediate-axes}
The RHS may have a different graph structure from the LHS, so its
intermediate tensors may retain different logical axes. To infer RHS axis
attributes, we first identify the axes that must survive at each node.
We then construct a retained-axis set $S_v$ using the operation's rules.
Axis ordering and attribute construction are handled in the following stages.

For a node $v$, let $\mathrm{In}_v$ be the set of inputs in the subtree rooted at $v$,
and let
$\sigma(a)$ be the inputs that contain logical axis $a$. The set $M_v$ of
\emph{mandatory axes} contains axes used in this subtree that are also needed in the
final output $\operatorname{Ord}_{\mathrm{out}}$ or in an input outside the subtree:
\begin{equation}
 M_v=\left\{a\in\Axes(L)\;\middle|\;
 \sigma(a)\cap \mathrm{In}_v\ne\varnothing
 \land (a\in\operatorname{set}(\operatorname{Ord}_{\mathrm{out}})\lor\sigma(a)\nsubseteq \mathrm{In}_v)
 \right\}.
 \label{eq:mandatory-axes}
\end{equation}
An operation may also have to keep an axis that is no longer needed.
For example, a contraction cannot sum out an axis present in only one
operand; it must keep that axis as a free axis.

Let $S_v$ be the axes that our algorithm keeps in the output of node
$v$.
This set includes every mandatory axis, so $M_v\subseteq S_v$.
Here $M_v$ specifies what must be preserved, while $S_v$ is the constructed
result. The inclusion alone does not determine $S_v$.
We call it the \emph{canonical live set} because the algorithm follows one
fixed rule: remove each unneeded axis at the first operation that can
legally remove it. We provide operation-specific rules for constructing
$S_v$ in Appendix~\ref{app:operation-construction-rules}.

For example, consider the intermediate node $v=\operatorname{dot}(X,Y)$ in
$\operatorname{dot}(\operatorname{dot}(X,Y),Z)$, with the LHS axis connections
\begin{equation}
 X:\vec B\Vert\vec I\Vert\vec K,\qquad
 Y:\vec B\Vert\vec K\Vert\vec J,\qquad
 Z:\vec B\Vert\vec J\Vert\vec L,\qquad
 \operatorname{Ord}_{\mathrm{out}}=\vec B\Vert\vec I\Vert\vec L.
\end{equation}
The intermediate retains $\vec B$ as batch axes, $\vec I$ for the final
output, and $\vec J$ to connect with $Z$.
The axes in $\vec K$ are needed only inside this subtree
and can be summed out. Thus, in this example,
\begin{equation}
 M_v=S_v=\operatorname{set}(\vec B)\cup\operatorname{set}(\vec I)
 \cup\operatorname{set}(\vec J).
\end{equation}
Other operations may retain additional axes, giving $M_v\subsetneq S_v$.

Applying these rules from the inputs upward keeps every mandatory axis,
and every retained axis comes from an input in the subtree. At the root,
the remaining axes must be exactly those required by the LHS output:
\begin{equation}
 S_{\mathrm{root}}=\operatorname{set}(\operatorname{Ord}_{\mathrm{out}}).
 \label{eq:root-live}
\end{equation}
Extra axes may survive intermediate nodes, but none may remain in the final
output. Appendix~\ref{app:mandatory-and-retained-axes-example} provides a concrete example
of how these axis sets are computed for batched matrix multiplication
followed by a row sum, including a case where $M_v\subsetneq S_v$.

\subsection{Why the Canonical Choice Does Not Lose Solutions}
\label{sec:canonical-choice-preserves-solutions}
We now show that choosing canonical intermediate axes does not lose valid
solutions.

Let $\operatorname{Ord}_v(\beta)$ be the ordered list of logical axes
in node $v$'s output when the RHS attributes are assigned by $\beta$.
The notation $S|_Q$ denotes the subsequence obtained by removing from $S$
every element not in $Q$, while preserving the order of the remaining elements. By
Lemma~\ref{lem:axis-exactness}, any valid RHS has the LHS partition, so its
intermediate axes can be named using LHS classes.

\begin{lemma}[Canonicalization preserves feasibility]
\label{lem:canonicalization}
Under Section~\ref{sec:formal-scope-and-guarantees}'s assumptions, with every input dimension
having extent at least two, if a valid assignment $\beta$ exists, a valid
assignment $\beta'$ exists with
\begin{equation}
 \forall v,\qquad
 \operatorname{Ord}_v(\beta')=\operatorname{Ord}_v(\beta)|_{S_v}.
 \label{eq:canonicalization}
\end{equation}
\end{lemma}
The lemma guarantees that whenever a valid assignment exists, one also
exists that keeps exactly $S_v$ at each node while preserving the relative
order of retained axes.

\paperparagraph{Example.}
Let $A,B$ be batched tensors with dimension orders $(b,i,k)$
and $(b,k,j)$, respectively. Suppose the source computes batched matrix
multiplication followed by a sum over batches:
\[
 T_{ij}=\sum_b(A_bB_b)_{ij}=\sum_{b,k}A_{bik}B_{bkj}.
\]
The same RHS graph, $\operatorname{sum}(\operatorname{dot}(A,B))$, can
compute this matrix under two different attribute assignments:
\begin{table}[!htbp]
\caption{Two RHS attribute assignments for batched matrix multiplication followed by a sum over batches.}
\label{tab:canonical-example}
\centering
\small
\begin{tabular}{@{}lllll@{}}\toprule
Assignment & Dot $(b_\ell,b_r,k_\ell,k_r)$ & $\operatorname{Ord}_v$ & $S_v$ & Sum axes\\\midrule
$\beta$ & $([0],[0],[2],[1])$ & $(b,i,j)$ & $\{i,j\}$ & $[0]$\\
$\beta'$ & $([],[],[0,2],[0,1])$ & $(i,j)$ & $\{i,j\}$ & $[]$\\\bottomrule
\end{tabular}
\end{table}
Here $v$ is the inner dot node; $b_\ell,b_r$ and $k_\ell,k_r$ list its
batch and contracting axes. Assignment $\beta$ contracts $k$ in the dot
and sums over $b$ afterward. The canonical assignment $\beta'$ contracts
both immediately, so the outer sum reduces no axes. Both produce the same
matrix, but $\beta'$ retains exactly $S_v=\{i,j\}$, satisfying
$\operatorname{Ord}_v(\beta')=\operatorname{Ord}_v(\beta)|_{S_v}$.

For TEPL, this means that synthesis can fix the intermediate axis sets
before determining their order and attributes, without losing feasible
rewrites within the defined fragment.

The proof is given in Appendix~\ref{app:canonicalization-preserves-feasibility}.

\subsection{Checking Axis Order}
\label{sec:checking-axis-order}
The canonical live sets determine which axes each RHS node retains.
We now check whether the RHS graph can arrange these axes in the order
required by the LHS output.

We propagate ordering requirements from the root toward the inputs.
Let $Req_v$ be an ordered list of axes that must appear as a subsequence
of node $v$'s output. Write $A\preceq B$ when $A$ is a subsequence of $B$.
This relation specifies relative order, not adjacency: $(i,x,j)$
satisfies the requirement $(i,j)$.
At the root, the requirement is the complete target order,
$Req_{\mathrm{root}}=\operatorname{Ord}_{\mathrm{out}}$.

At each node, we check whether the required order is compatible with
the operation's output order and pass the necessary requirements to its
children. For example, a contraction with left-free axis $i$ and
right-free axis $j$ cannot satisfy $(j,i)$ because left-free axes precede
right-free axes. A permutation above the contraction can satisfy this
requirement by reordering the result. It therefore passes an empty
requirement to the contraction. Requirements propagate to the inputs,
where they are checked against the fixed input orders.
Appendix~\ref{app:operation-construction-rules} gives the rules for each
operation.

Let $C_v$ denote the local ordering check at node $v$.
The check also requires $Req_v$ to contain distinct axes from $S_v$.
Combining all local checks with the root axis-set check gives
\begin{equation}
 \operatorname{check}(m)
 =[S_{\mathrm{root}}=\operatorname{set}(\operatorname{Ord}_{\mathrm{out}})]
 \land\bigwedge_{v\in V_R}C_v,
\label{eq:feasibility}
\end{equation}
where $V_R$ contains all RHS nodes, including inputs.
The root axis-set check ensures that exactly the target axes survive.
The ordering checks ensure that these axes can appear in the target order.
Once these checks pass, the next stage constructs the axis attributes.

\subsection{Constructing Axis Attributes}
\label{sec:constructing-axis-attributes}
Once $\operatorname{check}(m)$ succeeds, the remaining task is to turn
the selected axes and ordering requirements into RHS axis attributes.
The attribute builder $\operatorname{build}$ processes the RHS from inputs to root
in three steps.

\paperparagraph{Step 1: Construct an output order.}
At each node $v$, we construct an output order $\operatorname{Ord}_v$
that contains exactly the axes in $S_v$ and satisfies the required
relative order $Req_v$:
\begin{equation}
 \operatorname{set}(\operatorname{Ord}_v)=S_v,
 \qquad Req_v\preceq\operatorname{Ord}_v.
 \label{eq:order-invariant}
\end{equation}
The construction follows the child output orders and the operation's
rules under D2--D3. 

\paperparagraph{Step 2: Convert axes to positions.}
The notation $\operatorname{pos}(S,\vec G)$ lists the zero-based positions
of the axes in $\vec G$ within $S$, preserving the group order.
For example, if $\operatorname{Ord}_u=(i_1,k_1,i_2,k_2)$ and
$\vec K=(k_1,k_2)$, then
$\operatorname{pos}(\operatorname{Ord}_u,\vec K)=[1,3]$.
These positions form the axis-list attribute selecting $\vec K$.

\paperparagraph{Step 3: Combine node assignments.}
By D4, the independently constructed node assignments can be combined
into $\beta=\operatorname{build}(m)$. For a single component, this completes
the construction. For multiple terms, we combine the component checks and
assignments as in Appendix~\ref{app:extension-to-multiple-terms}.

The rewrite tree is fixed, but the
attribute lists may vary in length, so the same builder handles different
input ranks and shapes within the defined fragment. Both $\operatorname{check}$
and $\operatorname{build}$ operate on finite metadata lists.
This defines the mathematical functions; Section~\ref{sec:axis-mappings-to-attribute-builders}
describes TEPL's symbolic groups and expression generation.
